\chapter{Reti Neurali}

In questo capitolo viene presentata l'analisi e l'implementazione di reti neurali artificiali per la classificazione binaria del diabete. Data la natura del problema e lo sbilanciamento intrinseco del dataset, sono stati esplorati diversi approcci per migliorare le performance predittive, con particolare attenzione alla riduzione dei falsi negativi, cruciale in un contesto clinico-diagnostico.

\section{Architettura di Base}

\subsection{Struttura della Rete}

L'architettura adottata è un \textit{Multi-Layer Perceptron} (MLP) con la seguente configurazione:

\begin{itemize}
    \item \textbf{Layer di Input}: 21 neuroni (uno per ogni feature standardizzata)
    \item \textbf{Layer Nascosti}: tre strati con struttura decrescente (128 $\rightarrow$ 64 $\rightarrow$ 32 neuroni)
    \item \textbf{Layer di Output}: 1 neurone con attivazione sigmoidea
\end{itemize}

\subsubsection{Giustificazione della Profondità}

La scelta di tre layer nascosti rappresenta un compromesso tra:
\begin{enumerate}
    \item \textbf{Capacità rappresentativa}: una rete troppo superficiale potrebbe non catturare relazioni non lineari complesse tra le feature cliniche
    \item \textbf{Rischio di overfitting}: una rete troppo profonda tenderebbe a memorizzare i dati di training, specialmente con dataset sbilanciati
    \item \textbf{Efficienza computazionale}: il numero di parametri rimane gestibile mantenendo tempi di training ragionevoli
\end{enumerate}

\subsubsection{Giustificazione della Struttura Decrescente (128 $\rightarrow$ 64 $\rightarrow$ 32)}

La struttura \textit{a piramide} con numero decrescente di neuroni è motivata da:

\begin{itemize}
    \item \textbf{Estrazione gerarchica delle feature}: i primi layer (più ampi) apprendono rappresentazioni di basso livello delle feature originali, mentre i layer successivi (più stretti) combinano queste rappresentazioni in pattern di alto livello più astratti
    \item \textbf{Compressione dell'informazione}: la riduzione progressiva forza la rete a sintetizzare le informazioni più rilevanti per la classificazione
    \item \textbf{Regolarizzazione implicita}: il collo di bottiglia riduce la capacità della rete di memorizzare rumore nei dati
\end{itemize}

\subsection{Funzioni di Attivazione}

\subsubsection{ReLU nei Layer Nascosti}

Per i layer nascosti è stata scelta la funzione \textit{Rectified Linear Unit} (ReLU):

\begin{equation}
    f(x) = \max(0, x)
\end{equation}

Le motivazioni di questa scelta includono:

\begin{itemize}
    \item \textbf{Mitigazione del vanishing gradient}: a differenza delle funzioni sigmoidee o tangente iperbolica, ReLU non satura per valori positivi, permettendo gradienti non nulli durante la backpropagation
    \item \textbf{Sparsità delle attivazioni}: i neuroni con input negativo producono output nullo, creando rappresentazioni sparse che possono migliorare la generalizzazione
    \item \textbf{Efficienza computazionale}: il calcolo di ReLU richiede solo un confronto, risultando più veloce rispetto a funzioni esponenziali
    \item \textbf{Convergenza più rapida}: empiricamente, reti con ReLU convergono più velocemente durante l'addestramento
\end{itemize}

\subsubsection{Sigmoide nel Layer di Output}

Il layer di output utilizza la funzione sigmoide:

\begin{equation}
    \sigma(x) = \frac{1}{1 + e^{-x}}
\end{equation}

Questa scelta è obbligata per la classificazione binaria poiché:
\begin{itemize}
    \item Produce output nell'intervallo $[0, 1]$, interpretabile come probabilità di appartenenza alla classe positiva (diabete presente)
    \item È compatibile con la funzione di loss \textit{Binary Cross-Entropy}
    \item Permette di applicare threshold personalizzati per ottimizzare precision/recall
\end{itemize}

\subsection{Regolarizzazione con Dropout}

Tra ogni layer nascosto è stato inserito uno strato di \textbf{Dropout} con probabilità $p = 0.3$:

\begin{equation}
    h_i^{(l)} = \begin{cases}
        0 & \text{con probabilità } p \\
        \frac{h_i^{(l)}}{1-p} & \text{con probabilità } 1-p
    \end{cases}
\end{equation}

\subsubsection{Motivazioni del Dropout}

\begin{itemize}
    \item \textbf{Prevenzione dell'overfitting}: disattivando casualmente il 30\% dei neuroni durante ogni iterazione di training, si impedisce alla rete di sviluppare dipendenze eccessive da specifici neuroni
    \item \textbf{Ensemble implicito}: il dropout può essere interpretato come l'addestramento simultaneo di un ensemble di sotto-reti, migliorando la robustezza delle predizioni
    \item \textbf{Dataset sbilanciato}: con una classe dominante (86\% negativi), il dropout aiuta a prevenire che la rete impari pattern troppo specifici della classe maggioritaria
\end{itemize}

\subsubsection{Scelta del Valore $p = 0.3$}

Il valore $p = 0.3$ è stato scelto come compromesso:
\begin{itemize}
    \item Valori troppo bassi ($p < 0.2$) forniscono regolarizzazione insufficiente
    \item Valori troppo alti ($p > 0.5$) possono rallentare eccessivamente l'apprendimento e degradare le performance
    \item Il valore 0.3 è empiricamente efficace per reti di questa dimensione su problemi di classificazione tabellare
\end{itemize}

\section{Configurazione dell'Addestramento}

\subsection{Funzione di Loss: Binary Cross-Entropy}

Per l'addestramento è stata utilizzata la \textit{Binary Cross-Entropy} (BCE):

\begin{equation}
    \mathcal{L}_{BCE} = -\frac{1}{N} \sum_{i=1}^{N} \left[ y_i \log(\hat{y}_i) + (1 - y_i) \log(1 - \hat{y}_i) \right]
\end{equation}

dove $y_i \in \{0, 1\}$ è l'etichetta vera e $\hat{y}_i \in [0, 1]$ è la probabilità predetta.

\subsubsection{Giustificazione}

\begin{itemize}
    \item \textbf{Interpretazione probabilistica}: la BCE deriva dal principio di massima verosimiglianza per distribuzioni Bernoulli
    \item \textbf{Gradienti informativi}: penalizza maggiormente predizioni confidenti ma errate, fornendo segnali di errore più forti quando necessario
    \item \textbf{Compatibilità}: si integra naturalmente con l'output sigmoidale, evitando problemi numerici
\end{itemize}

\subsection{Ottimizzatore: Adam}

L'ottimizzazione dei pesi è affidata all'algoritmo \textbf{Adam} (\textit{Adaptive Moment Estimation}):

\begin{align}
    m_t &= \beta_1 m_{t-1} + (1 - \beta_1) g_t \\
    v_t &= \beta_2 v_{t-1} + (1 - \beta_2) g_t^2 \\
    \hat{m}_t &= \frac{m_t}{1 - \beta_1^t}, \quad \hat{v}_t = \frac{v_t}{1 - \beta_2^t} \\
    \theta_t &= \theta_{t-1} - \frac{\eta}{\sqrt{\hat{v}_t} + \epsilon} \hat{m}_t
\end{align}

\subsubsection{Vantaggi di Adam}

\begin{itemize}
    \item \textbf{Learning rate adattivo}: ogni parametro riceve un learning rate personalizzato basato sulla storia dei gradienti
    \item \textbf{Momentum}: la media mobile del primo momento ($m_t$) accelera la convergenza nelle direzioni consistenti
    \item \textbf{Normalizzazione dei gradienti}: la divisione per $\sqrt{\hat{v}_t}$ stabilizza l'addestramento in presenza di gradienti con scale diverse
    \item \textbf{Robustezza}: funziona bene con i valori di default ($\beta_1 = 0.9$, $\beta_2 = 0.999$, $\eta = 0.001$) senza necessità di tuning estensivo
\end{itemize}

\subsection{Iperparametri di Training}

\subsubsection{Numero di Epoche: 150}

Il numero di epoche è stato fissato a 150 per:
\begin{itemize}
    \item Permettere alla rete di convergere completamente
    \item Osservare eventuali fenomeni di overfitting nelle curve di validazione
    \item Bilanciare tempo di training e qualità del modello
\end{itemize}

\subsubsection{Batch Size: 256}

La dimensione del batch è stata impostata a 256 campioni:
\begin{itemize}
    \item \textbf{Stabilità del gradiente}: batch più grandi producono stime del gradiente meno rumorose
    \item \textbf{Efficienza GPU}: dimensioni potenze di 2 ottimizzano l'utilizzo della memoria
    \item \textbf{Generalizzazione}: batch non troppo grandi mantengono un certo livello di rumore nel gradiente, che può fungere da regolarizzatore implicito
\end{itemize}

\subsubsection{Validation Split: 10\%}

Il 10\% del training set è riservato alla validazione per:
\begin{itemize}
    \item Monitorare l'overfitting durante l'addestramento
    \item Valutare la generalizzazione senza toccare il test set
    \item Decidere eventuali strategie di early stopping
\end{itemize}

\section{Gestione dello Sbilanciamento delle Classi}

Il dataset presenta un forte sbilanciamento: circa l'86\% delle istanze appartiene alla classe negativa (non diabetici) e solo il 14\% alla classe positiva (diabetici). Questo sbilanciamento è problematico perché:

\begin{itemize}
    \item Un modello che predice sempre ``negativo'' otterrebbe un'accuracy dell'86\%
    \item In contesto clinico, i \textbf{falsi negativi} (mancata diagnosi) sono particolarmente gravi
    \item Le metriche tradizionali come l'accuracy sono fuorvianti
\end{itemize}

Sono state quindi esplorate tre strategie per mitigare questo problema.

\subsection{Approccio 1: Undersampling}

\subsubsection{Metodologia}

L'undersampling bilancia il dataset riducendo la classe maggioritaria:
\begin{enumerate}
    \item Si identificano tutti gli indici della classe minoritaria (positivi)
    \item Si campionano casualmente dalla classe maggioritaria lo stesso numero di istanze
    \item Si crea un nuovo dataset bilanciato (50\%-50\%)
\end{enumerate}

\subsubsection{Architettura Ridotta}

Per l'undersampling è stata adottata un'architettura più compatta (64 $\rightarrow$ 32 $\rightarrow$ 16 neuroni) con:
\begin{itemize}
    \item \textbf{Dropout aumentato al primo layer ($p = 0.5$)}: maggiore regolarizzazione per compensare il dataset ridotto
    \item \textbf{Meno neuroni}: dataset più piccolo richiede meno capacità per evitare overfitting
    \item \textbf{Meno epoche (100)}: convergenza più rapida su dataset bilanciato
    \item \textbf{Batch size ridotto (128)}: adattato alla dimensione ridotta del dataset
\end{itemize}

\subsubsection{Vantaggi e Svantaggi}

\textbf{Vantaggi}:
\begin{itemize}
    \item Training più veloce
    \item Classi bilanciate naturalmente
    \item Recall sulla classe positiva notevolmente migliorato
\end{itemize}

\textbf{Svantaggi}:
\begin{itemize}
    \item Perdita di informazione dalla classe maggioritaria
    \item Possibile perdita di pattern rilevanti
    \item Performance valutata su test set ridotto
\end{itemize}

\subsection{Approccio 2: Class Weights}

\subsubsection{Metodologia}

I class weights modificano la funzione di loss per penalizzare maggiormente gli errori sulla classe minoritaria:

\begin{equation}
    \mathcal{L}_{weighted} = -\frac{1}{N} \sum_{i=1}^{N} w_{y_i} \left[ y_i \log(\hat{y}_i) + (1 - y_i) \log(1 - \hat{y}_i) \right]
\end{equation}

I pesi sono calcolati come:
\begin{equation}
    w_c = \frac{N}{2 \cdot n_c}
\end{equation}

dove $N$ è il numero totale di campioni e $n_c$ è il numero di campioni della classe $c$.

\subsubsection{Configurazione}

Per il dataset in esame:
\begin{itemize}
    \item $w_0 \approx 0.58$ (classe negativa, maggioritaria)
    \item $w_1 \approx 3.57$ (classe positiva, minoritaria)
\end{itemize}

Questo significa che un errore sulla classe positiva costa circa 6 volte di più durante l'ottimizzazione.

\subsubsection{Vantaggi e Svantaggi}

\textbf{Vantaggi}:
\begin{itemize}
    \item Utilizza tutti i dati disponibili
    \item Non richiede modifica del dataset
    \item Implementazione semplice
\end{itemize}

\textbf{Svantaggi}:
\begin{itemize}
    \item Curve di validazione più rumorose
    \item Richiede tuning del bilanciamento
    \item Può causare oscillazioni durante l'ottimizzazione
\end{itemize}

\subsection{Approccio 3: Focal Loss}

\subsubsection{Formulazione}

La Focal Loss è una generalizzazione della Cross-Entropy che downweighta gli esempi ``facili'':

\begin{equation}
    \mathcal{L}_{FL} = -\alpha_t (1 - p_t)^\gamma \log(p_t)
\end{equation}

dove:
\begin{itemize}
    \item $p_t = \hat{y}$ se $y = 1$, altrimenti $p_t = 1 - \hat{y}$
    \item $\gamma$ è il \textbf{focusing parameter}
    \item $\alpha_t$ è il peso di bilanciamento delle classi
\end{itemize}

\subsubsection{Parametri Utilizzati}

\begin{itemize}
    \item \textbf{$\gamma = 2.0$}: valore standard che fornisce un buon bilanciamento. Valori più alti ($\gamma > 2$) focalizzano eccessivamente sugli esempi difficili, che potrebbero essere outlier o rumore
    \item \textbf{$\alpha$}: calcolato automaticamente dalla distribuzione delle classi:
    \begin{equation}
        \alpha = \frac{n_0}{n_0 + n_1} \approx 0.86
    \end{equation}
\end{itemize}

\subsubsection{Meccanismo di Funzionamento}

Il termine $(1 - p_t)^\gamma$ modula la loss:
\begin{itemize}
    \item Per esempi ben classificati ($p_t \rightarrow 1$): $(1 - p_t)^\gamma \rightarrow 0$, riducendo il loro contributo
    \item Per esempi difficili ($p_t \rightarrow 0$): $(1 - p_t)^\gamma \rightarrow 1$, mantenendo il loro peso
\end{itemize}

\subsubsection{Vantaggi e Svantaggi}

\textbf{Vantaggi}:
\begin{itemize}
    \item Focus automatico sugli esempi difficili
    \item Combina bilanciamento classi e difficoltà degli esempi
    \item Teoricamente più sofisticato
\end{itemize}

\textbf{Svantaggi}:
\begin{itemize}
    \item Maggiore onere computazionale
    \item Richiede tuning di parametri aggiuntivi ($\gamma$, $\alpha$)
    \item Nel caso in esame, miglioramenti marginali rispetto ai class weights
\end{itemize}

\section{Esperimenti con PCA}

\subsection{Motivazione}

L'analisi delle componenti principali (PCA) è stata applicata per verificare se la riduzione dimensionale potesse:
\begin{itemize}
    \item Rimuovere rumore e feature ridondanti
    \item Accelerare il training
    \item Migliorare la generalizzazione
\end{itemize}

\subsection{Risultati della PCA}

L'analisi ha rivelato che sono necessarie \textbf{14 componenti} su 21 feature originali per catturare l'80\% della varianza. Questo risultato è coerente con:
\begin{itemize}
    \item La bassa correlazione tra le feature (matrice di covarianza con valori off-diagonal piccoli)
    \item L'assenza di feature dominanti nella heatmap PCA
    \item L'informazione distribuita uniformemente tra le feature
\end{itemize}

\subsection{Architettura con PCA}

Per i dati trasformati tramite PCA è stata adottata un'architettura più compatta:
\begin{itemize}
    \item \textbf{Layer}: 32 $\rightarrow$ 16 $\rightarrow$ 8 neuroni
    \item \textbf{Input}: 14 componenti principali
    \item \textbf{Dropout}: mantenuto a 0.3
\end{itemize}

La riduzione della capacità è giustificata dalla dimensionalità ridotta dell'input e dalla natura decorrelata delle componenti principali.

\subsection{Conclusioni sulla PCA}

I risultati mostrano che la PCA non apporta miglioramenti significativi per questo problema:
\begin{itemize}
    \item Le performance sono comparabili o leggermente inferiori
    \item La compressione limitata (14/21 componenti) non giustifica la perdita di interpretabilità
    \item Le feature originali, essendo clinicamente significative, potrebbero contenere informazioni non catturate dalla varianza
\end{itemize}

\section{Metriche di Valutazione}

\subsection{Matrice di Confusione}

La matrice di confusione permette di analizzare:
\begin{itemize}
    \item \textbf{True Negatives (TN)}: correttamente identificati come non diabetici
    \item \textbf{True Positives (TP)}: correttamente identificati come diabetici
    \item \textbf{False Negatives (FN)}: diabetici non riconosciuti (errore critico)
    \item \textbf{False Positives (FP)}: falsi allarmi (errore meno grave)
\end{itemize}

\subsection{Curva ROC e AUC}

La curva ROC (\textit{Receiver Operating Characteristic}) visualizza il trade-off tra:
\begin{itemize}
    \item \textbf{True Positive Rate} (Sensibilità): $TPR = \frac{TP}{TP + FN}$
    \item \textbf{False Positive Rate}: $FPR = \frac{FP}{FP + TN}$
\end{itemize}

L'\textbf{Area Under Curve} (AUC) fornisce una misura aggregata:
\begin{itemize}
    \item AUC = 0.5: classificatore casuale
    \item AUC = 1.0: classificatore perfetto
    \item AUC $\approx$ 0.8: buona capacità discriminativa (ottenuta dai modelli)
\end{itemize}

\subsection{Classification Report}

Il report include:
\begin{itemize}
    \item \textbf{Precision}: $\frac{TP}{TP + FP}$ - affidabilità delle predizioni positive
    \item \textbf{Recall}: $\frac{TP}{TP + FN}$ - capacità di trovare tutti i positivi
    \item \textbf{F1-Score}: media armonica di precision e recall
\end{itemize}

\section{Analisi Comparativa dei Risultati}

\subsection{Modello Base (No Bilanciamento)}

Il modello senza correzioni per lo sbilanciamento mostra:
\begin{itemize}
    \item Alta accuracy globale (dominata dalla classe maggioritaria)
    \item Basso recall sulla classe positiva
    \item Elevato numero di falsi negativi
    \item Curva ROC con AUC $\approx$ 0.8 (buona separabilità delle classi)
\end{itemize}

\subsection{Undersampling}

L'approccio con undersampling produce:
\begin{itemize}
    \item Precision e recall bilanciati tra le classi
    \item Riduzione significativa dei falsi negativi
    \item Aumento dei falsi positivi (trade-off accettabile in contesto clinico)
    \item Curve di training/validation più stabili
\end{itemize}

\subsection{Class Weights}

L'utilizzo dei pesi bilanciati risulta in:
\begin{itemize}
    \item Recall sulla classe positiva comparabile all'undersampling ($\approx$ 0.78)
    \item Precision inferiore ($\approx$ 0.31)
    \item Curve di validazione più rumorose
    \item Utilizzo completo dei dati disponibili
\end{itemize}

\subsection{Focal Loss}

La Focal Loss mostra:
\begin{itemize}
    \item Risultati simili ai class weights
    \item Maggiore complessità computazionale
    \item Nessun miglioramento sostanziale nel caso in esame
\end{itemize}

\section{Considerazioni Cliniche}

In un contesto diagnostico per il diabete, è fondamentale considerare:

\subsection{Trade-off Falsi Negativi vs Falsi Positivi}

\begin{itemize}
    \item \textbf{Falso Negativo}: paziente diabetico non diagnosticato $\rightarrow$ conseguenze potenzialmente gravi (complicazioni non trattate)
    \item \textbf{Falso Positivo}: paziente sano classificato come diabetico $\rightarrow$ ulteriori test, ansia, ma conseguenze limitate
\end{itemize}

Pertanto, in questo contesto è preferibile minimizzare i falsi negativi anche a costo di più falsi positivi.

\subsection{Raccomandazioni d'Uso}

\begin{itemize}
    \item \textbf{Screening massivo}: preferire modelli con alto recall (undersampling o class weights)
    \item \textbf{Riduzione falsi allarmi}: preferire modello base o con PCA + class weights
    \item \textbf{Bilanciamento}: considerare threshold personalizzati sulla probabilità di output
\end{itemize}

\section{Conclusioni}

L'analisi condotta ha dimostrato che:

\begin{enumerate}
    \item L'architettura MLP con struttura decrescente (128 $\rightarrow$ 64 $\rightarrow$ 32) e dropout fornisce una base solida per la classificazione
    \item Lo sbilanciamento delle classi richiede necessariamente tecniche di mitigazione
    \item \textbf{Class Weights} e \textbf{Undersampling} producono risultati comparabili, con il primo che mantiene più informazione
    \item La \textbf{Focal Loss} non apporta miglioramenti significativi per questo specifico problema
    \item La \textbf{PCA} non è vantaggiosa data la bassa correlazione tra le feature
    \item La scelta del modello finale dovrebbe essere guidata dal contesto applicativo e dalla tolleranza ai diversi tipi di errore
\end{enumerate}
